% ==========================================================
% 每组人数有上下限的分配
% 难度：★★ 中档
% ==========================================================

\newcommand{\qProbPermMinMaxCapStem}{%
五名大学生前往观看冰球、速滑、花滑三场不同比赛。每场比赛至少有一名、至多有两名学生前往。甲同学不去观看冰球比赛。求分配方案数。
}

\newcommand{\qProbPermMinMaxCapAnswer}{%
$60$
}

\newcommand{\qProbPermMinMaxCapSolution}{%
总人数为 $5$，比赛有三场，每场人数只能为 $1$ 或 $2$。因此三场比赛的人数只能是 $(2,2,1)$。按照冰球比赛有一人还是两人分类。

\textbf{情形一：冰球比赛有一人}
该人不能是甲，所以从其余四人中选择：$4$ 种。剩余四人分配到速滑、花滑两场比赛，每场两人。为速滑比赛选择两人后，其余两人自动去花滑比赛：$\binom{4}{2}=6$ 种。该情形共有 $4\times 6=24$ 种。

\textbf{情形二：冰球比赛有两人}
甲不能去冰球，因此从其余四人中选择两人：$\binom{4}{2}=6$ 种。剩余三人分配到速滑、花滑两场比赛，人数分别为 $1$ 和 $2$。先选择哪一场只有一人，有 $2$ 种；再选择该名学生，有 $3$ 种。该情形共有 $6\times 2\times 3=36$ 种。

总方案数为 $24+36=60$ 种。
}

\newcommand{\qProbPermMinMaxCapTeachingNotes}{%
\textbf{【避坑与边界说明】}
\begin{itemize}
    \item 三场比赛互不相同，容量类型 $(2,2,1)$ 的不同排列需要区分。
    \item 当冰球比赛人数确定后，另外两场的容量结构随之确定。
    \item “剩余人员自动进入最后一组”时，不要重复计数。
\end{itemize}
}
